\chapter{Model Evaluation and Validation}
\label{ch:evaluation}

\epigraph{``A model is not good because it fits the training data well. A model is good because it predicts unseen data well.''}{}

\learninggoal{
	\item Design proper train--validation--test splits for reliable evaluation.
	\item Apply cross-validation and understand its bias--variance tradeoff.
	\item Diagnose overfitting using learning curves.
	\item Avoid common evaluation pitfalls: data leakage, selection bias, and multiple testing.
}

\section{The Evaluation Problem}

The central question in model evaluation: \emph{how will this model perform on new, unseen data?} We cannot answer this by testing on the data used to train the model, because the model has memorized it.

\begin{definition}[Generalization Error]
	The \textbf{generalization error} is the expected loss on new data drawn from the same distribution as the training data:
	\[
		\mathcal{R}(f) = \E_{(\mathbf{x}, y) \sim P}[L(y, f(\mathbf{x}))]
	\]
	We estimate this by holding out a test set.
\end{definition}

\section{Train--Validation--Test Split}

Always split data into three disjoint sets:

\begin{itemize}
	\item \textbf{Training set} (typically 60--80\%): used to fit model parameters.
	\item \textbf{Validation set} (10--20\%): used for hyperparameter tuning and model selection.
	\item \textbf{Test set} (10--20\%): used \emph{exactly once} to report final performance.
\end{itemize}

\begin{principle}[The Golden Rule of Testing]
	The test set may influence model selection \emph{zero times}. Once you evaluate on the test set, that number is final. Any iterative improvement based on test scores invalidates the estimate.
\end{principle}

\section{Cross-Validation}

When data is limited, a single validation split is unreliable. Cross-validation averages over multiple splits.

\subsection{K-Fold Cross-Validation}

\begin{enumerate}
	\item Split data into $K$ equal folds (typically $K = 5$ or $10$).
	\item For fold $k$: train on all folds except $k$, validate on fold $k$.
	\item Report the average performance across folds.
\end{enumerate}

\begin{lstlisting}[language=Python, caption=K-fold cross-validation]
from sklearn.model_selection import cross_val_score

scores = cross_val_score(model, X, y, cv=5,
                         scoring='neg_mean_squared_error')
print(f'MSE: {-scores.mean():.3f} +/- {scores.std():.3f}')
\end{lstlisting}

\subsection{Bias--Variance of Cross-Validation}

\begin{itemize}
	\item $K = 2$: low computational cost, high bias (training on only half the data).
	\item $K = n$ (leave-one-out, LOOCV): low bias, high variance, high computation.
	\item $K = 10$: standard choice; balances bias and variance.
\end{itemize}

\subsection{Stratified Cross-Validation}

For classification, preserve class proportions in each fold:

\[
	\frac{n_k^{\text{(fold)}}}{n^{\text{(fold)}}} \approx \frac{n_k}{n}
\]

\section{Learning Curves}

Plot training error and validation error as a function of training set size:

\begin{itemize}
	\item \textbf{High bias:} both curves converge to a high error; adding more data does not help.
	\item \textbf{High variance:} training error is low, validation error is high, and the gap persists as data increases.
\end{itemize}

\begin{principle}[Diagnosing with Learning Curves]
	If the validation curve plateaus well above the training curve, you have high bias (the model is too simple). If the curves are far apart and the training curve is near zero, you have high variance (the model is too complex).
\end{principle}

\section{Bootstrap and Confidence Intervals}

The bootstrap provides confidence intervals for performance metrics without parametric assumptions:

\begin{enumerate}
	\item Resample the dataset with replacement $B$ times (e.g., $B = 1000$).
	\item For each bootstrap sample, train and evaluate the model.
	\item The $2.5$th and $97.5$th percentiles of the $B$ scores form a $95\%$ confidence interval.
\end{enumerate}

\section{Hyperparameter Tuning}

Hyperparameters control model complexity and are set before training:

\begin{longtable}{p{3.5cm}p{4cm}p{4cm}}
	\toprule
	\textbf{Model}   & \textbf{Hyperparameters}                              & \textbf{Effect}             \\
	\midrule
	Ridge regression & $\alpha$                                              & Regularization strength     \\
	Random forest    & \texttt{n\_estimators}, \texttt{max\_depth}           & Ensemble size, tree depth   \\
	SVM              & $C$, $\gamma$                                         & Margin, RBF width           \\
	Neural network   & \texttt{lr}, \texttt{batch\_size}, \texttt{layers}    & Learning rate, architecture \\
	XGBoost          & \texttt{eta}, \texttt{max\_depth}, \texttt{subsample} & Learning rate, tree params  \\
	\bottomrule
\end{longtable}

\subsection{Grid Search}

Exhaustively evaluate all combinations of specified hyperparameter values. With $K$-fold CV, this requires $K \times \prod_j |\Theta_j|$ model fits.

\subsection{Random Search}

Sample hyperparameter values from distributions. More efficient than grid search when only a few hyperparameters matter~\cite{bergstra2012random}.

\begin{lstlisting}[language=Python, caption=Random search with CV]
from sklearn.model_selection import RandomizedSearchCV
from scipy.stats import uniform, randint

param_dist = {
    'n_estimators': randint(50, 500),
    'max_depth': randint(3, 20),
    'learning_rate': uniform(0.01, 0.3)
}
search = RandomizedSearchCV(model, param_dist,
                             n_iter=100, cv=5)
search.fit(X_train, y_train)
\end{lstlisting}

\section{Common Pitfalls}

\subsection{Data Leakage}

Information from the test set leaks into training:

\begin{itemize}
	\item Fitting scaler or imputer \emph{before} splitting.
	\item Using future data to predict the past (time series).
	\item Feature selection on the full dataset before splitting.
\end{itemize}

\begin{principle}[No Peeking]
	Any transformation that uses information from outside the training set is data leakage. Fit all transformations (scaling, imputation, PCA, feature selection) on the training set \emph{only}, then apply to validation/test sets.
\end{principle}

\subsection{Multiple Testing and p-Hacking}

Evaluating many models on the same test set inflates the chance of finding a spuriously good result. If you test 20 models, you expect one to achieve $p < 0.05$ by chance.

\subsection{Selection Bias}

If the training data is not representative of the deployment distribution, the evaluation is optimistic. Common causes: convenience sampling, temporal drift, non-response bias.

\section{Comparing Models Statistically}

To compare two models A and B:

\begin{itemize}
	\item \textbf{Paired t-test:} compute the per-fold difference $d_k = \text{error}_k^A - \text{error}_k^B$ and test $\bar{d} \neq 0$.
	\item \textbf{McNemar's test:} for classification, compare the number of points where A is correct and B is wrong, vs.\ vice versa.
	\item \textbf{Bayesian comparison:} use a Bayesian hierarchical model to estimate the probability that A is better than B.
\end{itemize}

\section{Exercises}

\begin{exercise}
	Take a dataset and perform a 70/15/15 train/validation/test split. Train a model, tune hyperparameters on the validation set, and evaluate on the test set once. Now suppose you repeat the tuning step 10 times, each time evaluating on the same test set. What happens to your estimate of generalization error?
\end{exercise}

\begin{exercise}
	Generate a learning curve for a decision tree on a synthetic dataset. Vary the training set size from 10 to 500. Plot training MSE and validation MSE. At what point does the validation error stop decreasing? What does this tell you about the value of additional data?
\end{exercise}

\begin{exercise}
	Design a cross-validation strategy for a time series forecasting task. Why is standard K-fold CV inappropriate? What alternative would you use?
\end{exercise}

\begin{exercise}
	You accidentally scaled the full dataset before splitting into train and test. How does this affect the test error estimate? Quantify the bias: is it optimistic or pessimistic?
\end{exercise}
